\documentclass[10pt,a4paper]{article}

\usepackage[T1]{fontenc}
\usepackage[utf8]{inputenc}
\usepackage[ngerman]{babel}
\usepackage{lmodern}
\usepackage{geometry}
\geometry{left=1.7cm,right=1.7cm,top=1.45cm,bottom=1.55cm}
\usepackage{setspace}
\setstretch{0.99}
\usepackage{enumitem}
\usepackage{tabularx}
\usepackage{array}
\usepackage{hyperref}
\hypersetup{colorlinks=true,linkcolor=blue,urlcolor=blue}
\setlist[itemize]{leftmargin=1.15em,itemsep=0.06em,topsep=0.08em,parsep=0pt,partopsep=0pt}
\setlist[enumerate]{leftmargin=1.3em,itemsep=0.06em,topsep=0.08em,parsep=0pt,partopsep=0pt}
\renewcommand{\arraystretch}{1.18}
\pagestyle{empty}

\begin{document}

\begin{center}
{\Large \textbf{Triagent: Kompaktes technisches Design-Dokument}}\\[0.15em]
{\large Zielbild, Programmablauf und aktueller Prototypstand für PIC GmbH}\\[0.35em]
\textbf{Ansprechpartner:} Alexander Hülsmann und Tassilo Bechtolsheim\\
\textbf{Adressat:} Dennis Itzel, PIC GmbH \quad
\textbf{Stand:} April 2026
\end{center}
\small

\vspace{0.25em}
\noindent\textbf{Einordnung}\\
Dieses Dokument ist eine kompakte Arbeitsversion, keine vollständige Spezifikation. Es beschreibt das geplante Betriebsmodell von Triagent und reflektiert danach den heutigen Prototypstand aus technischer Sicht. Ziel ist, Entwicklungsaufwand, Integrationspunkte und offene Risiken besser einschätzen zu können.

\section*{1. Zielbild: Lebenszyklus von Triagent in der Praxis}

\noindent\textbf{Produktrolle im Zielbild}\\
Triagent soll als Triage-Schicht zwischen gemeldeter verdächtiger E-Mail und finaler SOC- oder MSSP-Fallbearbeitung laufen. Das System soll kein Secure Email Gateway ersetzen, sondern nach der Meldung einer verdächtigen Mail strukturieren, vorsortieren, Evidenz sammeln und den Analysten beim Fallabschluss unterstützen.

\vspace{0.25em}
\noindent\textbf{Geplante Betriebsvarianten}
\begin{itemize}
  \item \textbf{On-prem / private deployment:} Betrieb in der Infrastruktur des Kunden oder in einer vom Kunden kontrollierten Umgebung. Wichtig für regulierte Umfelder und Teams mit Datenhoheitsanforderungen.
  \item \textbf{Cloud / managed deployment:} Betrieb als gehostete Variante für Teams, die weniger Betriebsaufwand möchten.
  \item \textbf{Hybrid-Option:} lokale Verarbeitung sensibler Artefakte mit optionalen externen Enrichment-Diensten, falls der Kunde diese freigibt.
\end{itemize}

\noindent\textbf{Schnittstellen, die im Zielbild eingerichtet werden müssten}\par\vspace{0.12em}
\noindent\begin{tabularx}{\textwidth}{@{}>{\bfseries}p{3.2cm}X@{}}
Mail-Intake & Outlook Add-in, manueller Upload von \texttt{.eml}/\texttt{.msg}, später optional Shared Mailbox, Microsoft Graph, IMAP oder API-basierte Weiterleitung.\\
Identität / Rollen & Lokale Benutzer und Rollen im Prototyp, später SSO / IdP, Gruppenmapping und Mandanten-/Kundenkontext für MSSPs.\\
Analyse-Engines & Aktivierbare Module für Header/Auth, URLs, Redirects, Anhänge, Lookalikes, Threat Intel, Sandbox-Anbindung und kundenspezifische Regeln.\\
Fallausgabe & PDF-/JSON-Fallbericht, IOC-CSV/JSON, später Ticketing- oder SIEM-Ausgabe, zum Beispiel Jira, ServiceNow, Splunk, Sentinel oder SOAR.\\
Audit / Compliance & Nachvollziehbarer Verlauf aus Intake, Analyse, Assist-Draft, Analystenentscheidung und Export.\\
\end{tabularx}

\vspace{0.35em}
\noindent\textbf{Geplanter Programmablauf als Aktivitätsdiagramm in Textform}
\begin{enumerate}
  \item Mitarbeiter meldet verdächtige E-Mail über Outlook Add-in oder das SOC lädt eine \texttt{.eml}/\texttt{.msg}-Datei hoch.
  \item Triagent speichert Originalnachricht und Anhänge unverändert als Artefakte.
  \item Parser normalisieren Header, Body, URLs, Anhänge, Absenderdaten und Message-Metadaten.
  \item Aktivierte Analyse-Engines erzeugen technische Signale, zum Beispiel SPF/DKIM/DMARC, URL-Redirects, Lookalikes, Attachment-Metadaten und Kontextindikatoren.
  \item Triage-Policy ordnet den Fall einer Queue zu, zum Beispiel \texttt{Needs investigation}, \texttt{Likely benign} oder \texttt{Uncertain}.
  \item Assist-Draft erzeugt einen Vorschlag für Disposition, Klassifikationscode, markierte Artefakte und Begründung.
  \item Analyst prüft Evidenz, passt Vorschlag bei Bedarf an und bestätigt oder verwirft die Resolution.
  \item Triagent schreibt Audit-Event, aktualisiert Fallstatus und erzeugt auf Wunsch PDF-/JSON-/IOC-Exporte.
  \item Später optional: Übergabe an Ticketing, SIEM, SOAR oder MSSP-Kundensystem.
\end{enumerate}

\noindent\textbf{Erweiterbarkeit der Erkennungs- und Triage-Engines}\\
Das Zielbild ist ein modulares Engine-Modell. Einzelne Engines sollen pro Deployment aktivierbar, deaktivierbar und konfigurierbar sein. Für Kunden wäre damit zum Beispiel steuerbar, ob externe URL-Auflösung, Sandbox-Abfragen, Threat-Intel-Lookups oder nur lokale Heuristiken erlaubt sind. Perspektivisch sollte jede Engine ein einheitliches Ergebnisformat liefern:
\begin{itemize}
  \item \texttt{engine\_id}, Version, Konfiguration
  \item Eingabeartefakte, zum Beispiel URL, Attachment, Header
  \item strukturierte Signale, Score-Beiträge und Begründungen
  \item Confidence, Fehlerstatus und Audit-Metadaten
\end{itemize}

\newpage

\section*{2. Technische Reflexion: aktueller Prototypstand}

\noindent\textbf{Architektur des aktuellen Prototyps}\\
\noindent\begin{tabularx}{\textwidth}{@{}>{\bfseries}p{3.4cm}X@{}}
Frontend & Next.js, React, TypeScript. Enthält Landing Page, Demo, Uploads, In-tray, Dashboard, Report Detail, Resolve Drawer, Settings und Admin-Sichten.\\
Backend & FastAPI in Python mit serviceorientierter Struktur für Parsing, Triage, Assist Drafts, Evidence Export, Audit, RBAC und Demo-Seeding.\\
Datenhaltung & PostgreSQL über SQLAlchemy und Alembic-Migrationen.\\
Objektspeicher & MinIO als S3-kompatibler Storage für Originalmails, Anhänge und Exporte.\\
Deployment & Docker Compose für lokalen und self-hosted Demo-Betrieb.\\
Client-Integration & Office.js Outlook Add-in mit API-Key-basierter Meldung an das Backend.\\
\end{tabularx}

\vspace{0.35em}
\noindent\textbf{Komponentendiagramm in Textform}\\
\texttt{Outlook Add-in / Web UI / API} \(\rightarrow\) \texttt{FastAPI Backend} \(\rightarrow\) \texttt{Parser + Analyse-Services} \(\rightarrow\) \texttt{Triage Policy + Assist Draft} \(\rightarrow\) \texttt{PostgreSQL + MinIO} \(\rightarrow\) \texttt{Report UI + Export + Audit}

\vspace{0.35em}
\noindent\textbf{Was heute bereits implementiert ist}\par\vspace{0.12em}
\noindent\begin{tabularx}{\textwidth}{@{}>{\bfseries}p{4.1cm}X@{}}
E-Mail-Intake & Upload und Parsing von \texttt{.eml} und \texttt{.msg}, Batch-Import synthetischer Demos, Outlook Add-in für gemeldete Mails.\\
Report-Workspace & Uploads, In-tray, Dashboard, Report Detail mit Details, Lookalikes, Authentication, URLs, Attachments, Transmission, X-Headers und Source.\\
Technische Analyse & SPF, DKIM, DMARC, ARC, URL-Extraktion, optionale Redirect-Auflösung, Same-org-Lookalike-Erkennung, Attachment-Metadaten, Raw Source Preservation.\\
Triage & Persistierte Triage-Buckets und Queue-Filter. Diese Logik ist bewusst als frühe Heuristik zu verstehen und noch nicht produktionsreif.\\
Assist Draft & Lokaler heuristischer Assist-Draft für Disposition, Klassifikation, Begründung und vorgeschlagene Artefakte. Aktuell Demo- und Prototypstand.\\
Resolution & Analyst-in-the-loop-Resolution mit Disposition, Klassifikationscode, Notiz und markierten Artefakten.\\
Exporte & Evidence Export als PDF, Markdown und JSON sowie IOC Export als CSV und JSON.\\
Governance & Session-RBAC, Rollen/Rechte, API Keys, Audit Events, Audit Export und Demo-User-Konzept.\\
Kampagnenbasis & Datenmodell und Backend-Mechaniken für Campaign Assignment, Recluster, Merge, Split und Campaign Evidence sind angelegt, aber noch nicht Produktfokus.\\
\end{tabularx}

\vspace{0.35em}
\noindent\textbf{Pseudo-Algorithmus der aktuellen Triage}
\begin{verbatim}
report = ingest(message)
preserve_original(report)
signals = parse_headers_urls_attachments(report)
signals += auth_summary(report)
signals += lookalike_detection(report)
signals += url_resolution(report)
bucket = triage_policy(signals, risk_score)
draft = assist_draft(signals, bucket)
analyst_reviews(report, signals, draft)
if analyst_confirms:
    write_resolution()
    write_audit_event()
    export_evidence_if_requested()
\end{verbatim}

\noindent\textbf{Bekannte technische Grenzen und Code-Risiken}
\begin{itemize}
  \item \textbf{Triage und Assist Draft sind noch frühe Heuristiken}. Sie funktionieren für die Demo- und synthetischen Fälle, brauchen aber reale Daten, Tuning und saubere Bewertungsmetriken.
  \item \textbf{Engine-Erweiterbarkeit ist noch nicht sauber abstrahiert}. Analysefunktionen existieren als Python-Services, aber noch nicht als stabiles Plugin- oder Engine-SDK.
  \item \textbf{Deployment ist Compose-first}. Für echten Betrieb fehlen noch Packaging, Secrets Management, Backup-/Restore-Automation, Observability und Update-Prozesse.
  \item \textbf{Integrationen sind noch begrenzt}. Outlook Add-in und Upload/API sind da, tiefe Ticketing-, SIEM-, SOAR- oder Sandbox-Integrationen sind noch offen.
  \item \textbf{SSO / Enterprise Auth ist nicht fertig}. Lokales RBAC existiert, LDAP-nahe Tests existieren, aber die vollständige Enterprise-IdP-Story ist offen.
  \item \textbf{Robustheit unter Last ist nicht validiert}. Es gibt lokale Builds und gezielte Backend-Tests, aber noch keine Lasttests, Penetrationstests oder formale Security Reviews.
\end{itemize}

\newpage

\section*{3. Meilensteine, Teststand und nächste Schritte}

\noindent\textbf{Fertig oder weitgehend vorhanden}
\begin{itemize}
  \item lauffähiger Demo-Stack mit Frontend, Backend, Datenbank und Objektspeicher
  \item E-Mail-Parsing, Originalartefakt-Speicherung und technische Report-Sichten
  \item einfache Queue- und Triage-Logik
  \item Assist-Draft und analyst-in-the-loop-Resolution
  \item PDF-/JSON-/IOC-Exporte
  \item Audit- und RBAC-Grundlagen
  \item public read-only Demo und reproduzierbares Demo-Reset
\end{itemize}

\noindent\textbf{Ausstehende Kernmeilensteine}
\begin{itemize}
  \item \textbf{M1: Engine-Konzept stabilisieren}. Einheitliches Interface, Aktivierung pro Deployment, Fehlerbehandlung, Versionierung und Audit-Kontext.
  \item \textbf{M2: Triage-Qualität validieren}. Reale oder anonymisierte Testdaten, Metriken für False Positives, False Negatives, Analystenzeit und Auto-Categorization-Rate.
  \item \textbf{M3: Integrationspfade priorisieren}. Outlook Add-in härten, danach Shared Mailbox / Graph oder Ticketing-Ausgabe entscheiden.
  \item \textbf{M4: Betriebshärtung}. Secrets, Backup, Monitoring, Logging, Update-Prozess, Mandantenfähigkeit und Deployment-Profile.
  \item \textbf{M5: Enterprise Auth}. SSO / IdP, Gruppenmapping, Rollenmodell und Audit-Anforderungen für Kundenumgebungen.
  \item \textbf{M6: Pilotfähigkeit}. Pilot-Runbook, Datenvereinbarung, Success Metrics, Support-Prozess und klare Grenzen der Automatisierung.
\end{itemize}

\noindent\textbf{Optionale spätere Meilensteine}
\begin{itemize}
  \item Kampagnen-Workflow als Hauptoberfläche für wiederkehrende Wellen
  \item Mandantenfähige MSSP-Sichten mit Kundenkontext
  \item Sandbox- und Threat-Intel-Adapter
  \item Case-Management- und SIEM-Integrationen
  \item Cloud-Angebot neben on-prem / private deployment
  \item Lern- und Feedbackschleife aus bestätigten Analystenentscheidungen
\end{itemize}

\noindent\textbf{Teststand}\\
Es existieren gezielte Backend-Tests für Parsing, Auth-Zusammenfassung, URL-Resolution, Lookalike-Erkennung, Triage, Assist-Drafts, Evidence Exports, synthetischen Korpus und LDAP-nahe Rollenzuordnung. Zusätzlich wird das Frontend regelmäßig über den Next.js-Build geprüft. Es gibt aber aktuell \textbf{keine belastbare Coverage-Zahl}, keine vollständigen UI-Regressionstests für alle Workflows und keine formalen Last- oder Security-Tests. Der Code ist damit \textbf{demo- und pilotfähig}, aber noch nicht als robustes Enterprise-Produkt einzustufen.

\vspace{0.35em}
\noindent\textbf{Einschätzung für PIC GmbH}\\
Der sinnvollste Beitrag eines technischen Partners läge aus unserer Sicht in drei Bereichen:
\begin{itemize}
  \item \textbf{Engineering:} Engine-Abstraktion, Deployment-Härtung, Integrationen und Qualitätssicherung.
  \item \textbf{Evaluierung:} Definition realistischer Testfälle, Messung der Triage-Qualität und Bewertung der Analysten-Workflows.
  \item \textbf{Vertrieb / Partnering:} Einschätzung, welche Betriebsvariante und welcher erste Integrationspfad für SOCs oder MSSPs am besten verkaufbar ist.
\end{itemize}

\noindent\textbf{Kurzfazit}\\
Triagent hat bereits einen funktionsfähigen Kern für reported-email triage, Evidence Review, Assist Drafts, Resolution und Exporte. Die größten offenen Arbeiten liegen nicht in einer weiteren Feature-Demo, sondern in Produktisierung, Engine-Erweiterbarkeit, Integrationen, Betriebshärtung und echter Validierung auf realistischen Workflows.

\end{document}
